\section{Introduction}\label{sec:intro}

\subsection{Motivation and main results}
Let $F$ be a non-archimedean local field of residue characteristic $p > 0$ and $G$ a connected reductive group over $F$. 
In their seminal work, Fargues and Scholze \cite{Geometrization} proposed a categorical version of the local Langlands correspondence in which the two sides of the correspondence appear in terms of the following categories:
\begin{enumerate}
    \item On the automorphic side, one considers $\mc{D}_{\lis}(\Bun_G)$ the category of $\Lambda \in \{\Flb, \Zlb, \Qlb\}$ étale sheaves on the stack of $G$-torsors on the Fargues--Fontaine curve. 
    \item On the Galois side, one considers the category of coherent sheaves on $\Par_G$ the stack of $L$-parameters, introduced independently in \cite{Geometrization}, \cite{ZhuCoherentSheaves} and \cite{DHKM}. 
\end{enumerate}
The main result of \cite{Geometrization} is the construction of a spectral action. 
\begin{thm}[\protect{\cite[Theorem X.0.1]{Geometrization}}]\label{thm:intro-FS-spectral-action}
    Assuming that $\ell$ does not divide the order of $\pi_{1}(\widehat{G})_{\tn{tor}}$ if $\Lambda \in \{\Flb, \Zlb\}$, there exists a canonically defined action, called the spectral action, of $\Perf(\Par_G)$ on $\mc{D}_{\lis}(\Bun_G)$. 
\end{thm}

Once this action is in place, one can formulate the following categorical local Langlands correspondence. 
\begin{conj}[\protect{\cite[Conjecture X.3.5]{Geometrization}}]\label{conj:geometrization-intro}
    Assume $G$ is quasi-split and $\ell$ as in Theorem \ref{thm:intro-FS-spectral-action}. For $(U,\psi)$ a fixed Whittaker datum for $G$, there exists a canonical $\Perf(\Par_G)$-linear and Whittaker normalized equivalence 
    \begin{equation}\label{eq:conj-geometrization-intro}
        \mc{D}_{\lis}(\Bun_G) \cong \IndCoh_{\nilp}(\Par_G).
    \end{equation}
\end{conj}

The Whittaker normalization in the statement of Conjecture \ref{conj:geometrization-intro} means that the equivalence \eqref{eq:conj-geometrization-intro} sends the Whittaker sheaf to the structure sheaf of $\Par_G$. This normalization is also the reason why it is necessary to assume that $G$ is quasi-split.

From the point of view of the theory of representations of $p$-adic groups, this categorical local Langlands conjecture misses certain inner forms of $G$. Specifically, the stack $\Bun_G$ only witnesses extended pure inner forms of $G$. One critical example of this phenomenon is $\SL_n$, none of whose non-trivial inner forms are pure. 

This failure of the categorical conjecture to capture certain inner forms is not a new problem; the statement of the refined local Langlands correspondence from \cite{Vog93} requires a Whittaker normalization and thus only applies in full generality to pure inner forms of a quasi-split group. 

To access all inner forms in a manner allowing for endoscopy, Kaletha \cite{Kal16} introduced the notion of rigid inner forms of $G$ and subsequently proposed in \cite[\S 5.4]{Kal16} an expected extension of the refined local Langlands correspondence which incorporates all inner forms of $G$ and allows for a general statement of the endoscopic character identities. In more recent work, Fargues \cite{Fargues22} has proposed a geometric construction to explain how to frame rigid inner forms within the geometrization program. Specifically, one defines a certain gerbe\footnote{We use the notation $\mathfrak{X}_{S}$ for this gerbe only in the introduction for ease of exposition; in the main body of the paper, this gerbe is denoted by $[\mc{T}_{S}/\tilde{t}_{S}]$.} $\mathfrak{X}_{S}$ over the Fargues--Fontaine curve $X_S$ (for some test perfectoid space $S$) banded by $u := \varprojlim_{E/F,n} \mathrm{Res}_{E/F}(\mu_{n})$ and defines 
$$\Bun_G^e$$
the `extended stack of $G$-torsors' as the stack of $G$-torsors on the gerbe $\mathfrak{X}_{S}$ which are basic (the transition maps for $u$ are given by norms and $n/m$ power maps for $m \mid n$; see Definition \ref{defi:bandmap} for what we mean by ``basic'' here).  

\begin{thm}[\cite{Fargues22}]
    The stack $\Bun_G^e$ is a small $v$-stack and decomposes canonically as 
    \begin{equation*}
        \Bun_G^e \cong \bigsqcup_{\lambda \in \Hom_{F}(u,Z_{G})} \Bun_{G}^{e,\lambda}.
    \end{equation*}
   
    Moreover, there is a further non-canonical decomposition
    \begin{equation} \label{eq:decomp-bun-g-e-intro}
        \bigsqcup_{\lambda \in \Hom_{F}(u,Z_{G})} \Bun_{G}^{e,\lambda}  \cong  \bigsqcup_{\lambda \in \Hom_{F}(u,Z_{G})} \Bun_{G_{b_{\lambda}}},
    \end{equation}
     induced by twisting isomorphisms $\Bun_{G}^{e,\lambda} \to  \Bun_{G_{b_{\lambda}}}$ for all $\lambda$, where the collection of groups $\{G_{b_{\lambda}}\}_{\lambda}$ represents all rigid inner twists of $G$ up to extended pure inner twists. 
\end{thm}
In very much the same way that open strata of $\Bun_G$ are classifying stacks of the form $\pt/\underline{G_b(F)}$ for extended pure inner forms of $G$, open strata of $\Bun_G^e$ are classifying stacks of rigid inner twists of $G$.

On the coherent side, one introduces \cite[\S 12.1]{Fargues22} (originally appearing in \cite{Vog93} and \cite{Kaletha18}) the extended dual group $\widehat{G}^e$ defined as 
\begin{equation*}
    \widehat{G}^e = \varprojlim_{Z} \widehat{G/Z}
\end{equation*}
where the limit runs through the set of all finite central subgroups of $G$.
\begin{rque}
    If $G$ is semisimple, we have $\widehat{G}^e = \widehat{G}_{\tn{sc}}$ the simply connected cover of $\widehat{G}$.
\end{rque}
Following \cite{Kal16} and \cite{Fargues22}, one defines the extended stack of parameters. 
\begin{defi}
    Let $\Par_G^{\square}$ be the stack of framed $L$-parameters (also denoted by $Z^1(W_F, \widehat{G})$ in \cite{DHKM}), so that $\Par_G = \Par_G^{\square}/\widehat{G}$. The extended stack of parameters is 
    $$\Par_G^e = \Par_G^{\square}/\widehat{G}^e,$$
    where the $\widehat{G}^{e}$-action is inflated along $\widehat{G}^{e} \to \widehat{G}$. 
\end{defi}

It follows from Theorem \ref{thm:intro-FS-spectral-action} and the decomposition \eqref{eq:decomp-bun-g-e-intro} that the category $\mc{D}_{\lis}(\Bun_G^e)$ can be equipped with an action of $\Perf(\Par_G)$ (provided $\ell$ is not too small for modular coefficients). Our main theorem is as follows, which resolves part (1) of \cite[Conjecture 12.8]{Fargues22}: 
\begin{thm}[Theorem \ref{thm:extended-spectral-action-on-bun_G}]\label{thm:intro-extended-spectral-action}
    Assume $\ell$ is as in Theorem \ref{thm:intro-FS-spectral-action}. There is a canonically defined action of $\Perf(\Par_G^e)$ on $\mc{D}_{\lis}(\Bun_G^e)$ which extends the Fargues--Scholze spectral action on each of the components coming from the decomposition \eqref{eq:decomp-bun-g-e-intro}.
\end{thm}

With this action in place, we can formulate an extended version of the categorical conjecture.
\begin{conj}[\protect{\cite[Conjecture 12.8.2]{Fargues22}}, Conjecture \ref{conj:extended-categorical-equivalence}]\label{conj:intro-extended-cat-equiv}
    Assume that $\ell$ is as in Theorem \ref{thm:intro-FS-spectral-action} and $G$ is quasi-split. 
    There exists a canonical $\Perf(\Par_G^e)$-linear and Whittaker normalized equivalence of categories
    $$\mc{D}_{\lis}(\Bun_G^e) \cong \IndCoh_{\nilp}(\Par_G^e).$$
\end{conj}

Following results of \cite{Zou24}, we prove the conjecture when $G$ is a torus. 
\begin{thm}[Theorem \ref{thm:case-of-tori}]
    Assume $G$ is a torus, then Conjecture \ref{conj:intro-extended-cat-equiv} holds for $G$.
\end{thm}

We also expect that it should be possible to reduce Conjecture \ref{conj:intro-extended-cat-equiv} to Conjecture \ref{conj:geometrization-intro}, this takes the form of the following statement. 
\begin{conj}[Conjecture \ref{conj:reduction-extended-to-classical}]\label{conj:intro-reduction-to-classical}
    The natural functor 
    \begin{equation}\label{eq:reduction-to-classical-intro}
        \mc{D}_{\lis}(\Bun_G) \otimes_{\Ind\Perf(\Par_G)} \Ind\Perf(\Par_G^e) \to \mc{D}_{\lis}(\Bun_G^e)
    \end{equation}
    is an equivalence. 
\end{conj}

In particular, if Conjectures \ref{conj:intro-reduction-to-classical} and \ref{conj:geometrization-intro} hold then Conjecture \ref{conj:intro-extended-cat-equiv} holds. 

\begin{thm}[Theorem \ref{thm:reduction-connected-center-case}]\label{thm:intro-reduction-to-classical}
    If $Z_G$ is connected and $H^1(F,Z_G) = 0$ then Conjecture \ref{conj:intro-reduction-to-classical} holds. 
\end{thm}

There is also the recent result of Hansen--Mann:
\begin{thm}[\cite{HansenMann26}]
    Let $G = \mathrm{GL}_{n}$, $\Lambda = \ov{\mathbb{Q}_{\ell}}$, and $F$ a finite extension of $\mathbb{Q}_{p}$. Assuming \cite[Conjecture 1.6.2]{HansenMann26} (compatibility with Eisenstein functors), the spectral action induces a Whittaker-normalized natural equivalence 
    \begin{equation*}
        \mc{D}_{\tn{lis}}(\Bun_{\GL_{n}}) \xrightarrow{\sim} \IndCoh(\Par_{\GL_{n}}).
    \end{equation*}
\end{thm}

Combining the above theorem with  Theorem \ref{thm:intro-reduction-to-classical} implies that we can prove the case of $\GL_n$.

\begin{corol}
    Assume $F$ is a finite extension of $\mathbb{Q}_p$, $G = \GL_n$, $\Lambda = \Qlb$ and \cite[Conjecture 1.6.2]{HansenMann26} holds, then Conjecture \ref{conj:intro-extended-cat-equiv} holds. 
\end{corol}

\subsection{Outline of the construction}

Let us outline the key steps in the proof of Theorem \ref{thm:intro-extended-spectral-action}, assuming now that the ring of coefficients $\Lambda$ is a $\Z_{\ell}[\sqrt{q}]$-algebra. First note that the stack $\Par_G^e$ fits in the following Cartesian diagram 
\[\begin{tikzcd}
	{\Par_G^e} & {\pt/\widehat{G}^e} \\
	{\Par_G} & {\pt/\widehat{G}}
	\arrow[from=1-1, to=1-2]
	\arrow[from=1-1, to=2-1]
	\arrow[from=1-2, to=2-2]
	\arrow[from=2-1, to=2-2]
\end{tikzcd}\]

\begin{thm}[Theorem \ref{thm:relative-tensor-product}]
    Assume $\ell$ is as in Theorem \ref{thm:intro-FS-spectral-action}.
    The above diagram induces an equivalence of categories 
    $$\Ind\Perf(\Par_G^e) = \Ind\Perf(\Par_G) \otimes_{\Ind\Perf(\pt/\widehat{G})} \Ind\Perf(\pt/\widehat{G}^e).$$
\end{thm}

Consequently, the main piece of structure required to extend the action of $\Perf(\Par_G)$ is the data of an action of $\Perf(\pt/\widehat{G}^{e})$ extending the one of $\Perf(\pt/\widehat{G})$. Let us briefly recall its construction from \cite{Geometrization}. We fix once and for all a geometric point 
$$\Spa(C) \xrightarrow{y} \Div^{1}:=\Div^1_F$$
where $C$ is an algebraically closed perfectoid field of characteristic $p > 0$. The action of $\Perf(\pt/\widehat{G})$ on $\mc{D}_{\lis}(\Bun_G)$ is given by Hecke functors, i.e. let $V \in \Rep_{\Lambda}\widehat{G}$ and let $\mc{S}_{V}$ denote the corresponding Satake sheaf on $\Hk^{\loc}_G$, the local Hecke stack of $G$ (over $\Spd(C)$). There is a diagram 
\[\begin{tikzcd}
	{\Hk^{\loc}_G} & {\Hk_G} & \\
	{\Bun_G} && {\Bun_G \times \Spd(C)}
	\arrow["\varepsilon"{description}, from=1-2, to=1-1]
	\arrow["{\overleftarrow{h}}"{description}, from=1-2, to=2-1]
	\arrow["{\overrightarrow{h}}"{description}, from=1-2, to=2-3]
\end{tikzcd}\]
where $\Hk_G$ is the moduli stack of modifications of $G$ torsors at the point $y$. The corresponding endofunctor of $\mc{D}_{\lis}(\Bun_G)$, when $\Lambda$ is a torsion ring is given by 
\begin{equation}\label{eq:hecke-operator-intro}
    A \mapsto \overrightarrow{h}_{!}(\overleftarrow{h}^*(A) \otimes \varepsilon^*(\mc{S}_V)),
\end{equation}
upon identifying $\mc{D}_{\lis}(\Bun_G) \cong \mc{D}_{\lis}(\Bun_G \times \Spd(C))$. 

In general, for $\ell$-adic coefficients, one needs to use the formalism of solid sheaves as in \cite[Section IX]{Geometrization}.

Passing to the extended setup, the main technical construction we provide is that of an extended Hecke stack $\Hk^e_G$ and its local version $\Hk^{e,\loc}_{G}$ fitting into a diagram 
\[\begin{tikzcd}
	{\Hk^{e,\loc}_G} & {\Hk_G^e} & \\
	{\Bun_G^e} && {\Bun_G^e \times \Spd(C).}
	\arrow["\varepsilon"{description}, from=1-2, to=1-1]
	\arrow["{\overleftarrow{h}}"{description}, from=1-2, to=2-1]
	\arrow["{\overrightarrow{h}}"{description}, from=1-2, to=2-3]
\end{tikzcd}\]
These stacks are defined in Section \ref{sec:extendedheckedef}. 
There are two key subtleties in the definition and construction of these Hecke stacks.

The first one is that the stack $\Hk_{G}^e$ is not naively the stack of tuples $(x, \mathcal{P}_0, \mathcal{P}_1, \phi)$ where $x : S \to \Div^1$ is a divisor on $X_S$, $\mathcal{P}_i$ are $G$-torsors on $\mathfrak{X}_{S}$ and $\phi$ is an isomorphism $\mathcal{P}_0 \cong \mathcal{P}_1$ away from the preimage of the Cartier divisor corresponding to $x$ in $\mathfrak{X}_{S}$. This definition yields a stack which we call $\Hk_{G}^{e, \tn{naive}}$ (resp. $\Hk_{G}^{e,\loc, \tn{naive}}$ for the local variant) defined in Section \ref{sec:hecke-stacks}. There is a natural inclusion 
$$\Hk_{G}^{e, \tn{naive}} \subset \Hk_{G}^{e},$$
which (after base-changing to $\Spd(C)$) is a union of connected components of $\Hk_{G}^e$. From the point of view of geometric Satake, the stack $\Hk_{G}^{e, \tn{naive}}$ only allows for modifications of torsors indexed by $\Rep(\widehat{G})$ and does not witness $\Rep(\widehat{G}^e)$. Correspondingly, the `naive' Hecke operators that one obtains in this way are the ones that yield the action of $\Perf(\Par_G)$ on $\mc{D}_{\lis}(\Bun_G^e)$. 

To enlarge this Hecke stack, we forcibly introduce some twists, see Definitions \ref{defi:elochck} and \ref{defi:eglobhck}, by twisting modifications between $G$-torsors by certain $Z_G$-torsors which exist only on the punctured curve (or punctured disks in the local case). 

The resulting effect of introducing these twists is that one, \emph{a priori}, loses the convolution structure of the local Hecke stacks. This is the second subtlety of the construction: We have to also redefine a convolution product to equip the category of sheaves on $\Hk_{G}^{e, \loc}$ with a monoidal structure. This is done in Section \ref{sec:Hckconv}. Doing convolution in this way does not glue over $\Div^1$ and the corresponding Hecke operator only exists over $\Spd(C)$---it turns out that this is an expected feature of the construction, see Remark \ref{rque:non-gluing-to-div-1}. 

Once the convolution is suitably constructed, we can move on to proving the relevant Geometric Satake equivalence. 
\begin{defi}[Definition \ref{defi:satake-category}]
    Let $\Sat_G^e \subset \mc{D}(\Hk^{e, \loc}_G)$ denote the full subcategory of perverse sheaves on $\Hk^{e, \loc}_G$ which are $\tn{ULA}$ over $\Spd(C)$ and $\Lambda$-flat. 
\end{defi}

\begin{thm}[Extended Geometric Satake, Theorem \ref{thm:satake}]
    There is an equivalence of monoidal categories
    $$\Sat_G^e \cong \Coh^{\heartsuit, \Lambda-\tn{flat}}(\pt/\widehat{G}^e \times_{\pt/\widehat{G}} \pt/\widehat{G}^e)$$
    where the right-hand side is equipped with the convolution monoidal structure. 
\end{thm}

Consider the relative diagonal 
$$\Delta : \pt/\widehat{G}^e \to \pt/\widehat{G}^e \times_{\pt/\widehat{G}} \pt/\widehat{G}^e.$$
The functor $\Delta_*$ provides a monoidal functor $\Rep_{\Lambda}\widehat{G}^e \to \Sat_G^e$ which, in turn, yields the desired Hecke operators using the same formula as \eqref{eq:hecke-operator-intro}. We refer to Section \ref{sec:Extended-Hecke-Action} for the details. 

\subsection{Relation with \cite{GLCV}}
In the proof of the Geometric Langlands conjectures, in the final steps of the proof, the authors do a sequence of reductions: from a general reductive group $G$ to $G$ semisimple and then to $G$ semisimple and simply connected. To do this second reduction, they use a $2$-categorical Fourier--Mukai transform statement and obtain in particular variants of the Geometric Langlands equivalence for unramified inner forms \cite[Corollary 8.10.2]{GLCV} of the group $G$ over the algebraic curve $X$ at play in \emph{loc. cit.} As the formalism of rigid inner forms has been constructed to handle inner forms of $G$, it is reasonable to expect that a variant of the construction of \cite{GLCV} in the Fargues--Scholze context would provide categorical equivalences for inner forms of $G$. The goal of this section is to have a somewhat informal discussion about the relationship between Theorem \ref{thm:intro-extended-spectral-action} and the construction of \cite{GLCV}. We will assume $\Lambda$ is torsion and $G$ is semisimple for simplicity, setting $Z:= Z_{G}$ temporarily. We also assume that $F$ has characteristic zero.

\begin{defi}
    Let $\Gamma$ be a finite flat abelian group scheme over $F$, we define $\Ge(\Gamma)$ to be the stack of gerbes banded by $\Gamma$, defined as 
    \begin{align*}
        S \in \Perfd \mapsto \{\tn{\'{e}tale gerbes over} \ X_S\}.
    \end{align*}
\end{defi}

Dually, we will denote by $\Ge^{\alg}(\pi_1(\widehat{G})(1))$ the condensed algebraic stack over $\Lambda$ given by 
$$\Ge^{\alg}(\pi_1(\widehat{G})(1)) := \Maps(B\Weil_F, B^2(\pi_1(\widehat{G})(1))).$$
We will now make the following set of assumptions about sheaves of categories on these stacks : 
\begin{enumerate}
    \item the stack $\Ge^{\alg}(\pi_1(\widehat{G})(1))$ has trivial condensed structure and is determined by a single $\infty$-groupoid,
    \item the stack $\Ge(Z)$ is discrete, i.e. as stacks over $\Perfd$, they are pulled back from some $\infty$-groupoid, and that there is an equivalence of categories between sheaves of categories on $\Ge(Z)$ and $\Ge^{\alg}(Z)$.
    \item There is a well defined $2$-categorical Fourier--Mukai duality of algebraic stacks over $\Lambda$ (see \cite[Section 8.1]{GLCV} for a definition) 
    \begin{equation}\label{eq:two-fm-transform}
        \Ge^{\alg}(\pi_1(\widehat{G}(1))) \times \Ge^{\alg}(Z) \to B^2\Gm.
    \end{equation}
\end{enumerate}
Note that all of these assumptions are reasonable; the last two are analogues of the descriptions of the corresponding stacks in \cite[Sections 8.3 and 8.4]{GLCV} and the first one should follow from the finiteness of $\pi_1(\widehat{G})(1)$.

Following \cite{GLCV}, we introduce two sheaves of categories : 
\begin{enumerate}
    \item consider the map $\pi^{\autom} : \Bun_{G_{\ad}} \to \Ge(Z)$, then define, as in \cite[\S 8.6.3]{GLCV}
    $$\mc{D}^{\autom} := \pi^{\autom}_*\mc{D}_{\et,\Bun_{G_{\ad}}}$$
    where $\mc{D}_{\et,\Bun_{G_{\ad}}}$ is the sheaf of categories whose global sections are $\mc{D}_{\et}(\Bun_{G_{\tn{ad}}})$. We will view it as a sheaf of categories on $\Ge^{\alg}(Z)$, by the above assumption.
    \item As the stack $\Par_G$ is $1$-affine, the spectral action of $\Perf(\Par_G)$ on $\mc{D}(\Bun_G)$ determines a sheaf of categories $\mc{D}_{\Par_G}$ on $\Par_G$\footnote{For the purposes of this informal discussion, we ignore the differences between $\QCoh$ and $\Ind\Perf$.}.  Consider the map $\pi^{\spec} : \Par_G \to \Ge^{\alg}(\pi_1(\widehat{G})(1))$, then as in \cite[\S 8.6.6]{GLCV}, one defines 
    $$\mc{D}^{\spec} := \pi^{\spec}_*\mc{D}_{\Par_G}.$$
\end{enumerate}

\begin{conj}[Analogue of \protect{\cite[Theorem 8.6.8]{GLCV}}]\label{conj:two-categorical-fourier-mukai-transform}
    Under the $2$-categorical Fourier--Mukai transform \eqref{eq:two-fm-transform}, we have 
    \begin{equation}
        2\tn{-FM}(\mc{D}^{\spec}) = \mc{D}^{\autom},
    \end{equation}
    up to a sign twist. 
\end{conj}

Let us now relate Conjecture \ref{conj:two-categorical-fourier-mukai-transform} to Theorem \ref{thm:intro-extended-spectral-action}. Recall that we denote by $u$ the group banding the gerbe $\Kal$ and that there is a canonical isomorphism 
$$\Hom_F(u, Z) = X^*(\pi_1(\widehat{G})(1)),$$
and that we have 
$$\Bun_G^e = \bigsqcup_{\lambda \in \Hom_F(u, Z)} \Bun_G^{e, \lambda}.$$
In particular, there is a Cartesian diagram 
\[\begin{tikzcd}
	{\Bun_{G_{\ad}}} & {\Bun_{G}^e} \\
	{\Ge(Z)} & {\Hom_F(u,Z)}
	\arrow["{\pi^{\autom}}"', from=1-1, to=2-1]
	\arrow[from=1-2, to=1-1]
	\arrow["{\pi^e}", from=1-2, to=2-2]
	\arrow["{\ev_{\Kal}}", from=2-2, to=2-1]
\end{tikzcd}\]
Dually, there is a diagram  
\[\begin{tikzcd}
	{\Par_G^e} & {\Par_G} & \\
	{B\widehat{G}^e} & {B\widehat{G}} & {\Ge^{\alg}(\pi_1(\widehat{G})(1))} \\
	\pt & {B^2\pi_1(\widehat{G})(1)}
	\arrow[from=1-1, to=1-2]
	\arrow[from=1-1, to=2-1]
	\arrow[from=1-2, to=2-2]
	\arrow[from=1-2, to=2-3]
	\arrow[from=2-1, to=2-2]
	\arrow[from=2-1, to=3-1]
	\arrow[from=2-2, to=3-2]
	\arrow["\ev"{description}, from=2-3, to=3-2]
	\arrow["q"', from=3-1, to=3-2]
\end{tikzcd}\]
where $\ev : \Ge^{\alg}(\pi_1(\widehat{G})(1)) \to B^2\pi_1(\widehat{G})(1)$ is induced by the map $\pt \to B\Weil_F$. Note that this map is $1$-affine. Assuming Conjecture \ref{conj:two-categorical-fourier-mukai-transform} and the compatibilities of $2$-categorical Fourier-Mukai transforms established in \cite[Section 8.4]{GLCV}, we obtain an equivalence
\begin{equation}\label{eq:reduction-intro}
    q^*\ev_{\pt,*}\mc{D}^{\spec} \cong \ev_{\Kal}^*\mc{D}^{\autom}.
\end{equation}

Unwinding the definitions of those terms, the left-hand side is 
$$q^*\ev_{\pt,*}\mc{D}^{\spec} = \mc{D}(\Bun_G) \otimes_{\Ind\Perf(\pt/\widehat{G})} \Ind\Perf(\pt/\widehat{G}^e),$$
and the right-hand side is 
$$\ev_{\Kal}^*\mc{D}^{\autom} = \mc{D}(\Bun_G^e).$$
It follows that \eqref{eq:reduction-intro} amounts to a version of Conjecture \ref{conj:intro-reduction-to-classical}. Note however that it is not clear, even assuming Conjecture \ref{conj:two-categorical-fourier-mukai-transform}, that Conjecture \ref{conj:intro-reduction-to-classical} holds as it is not clear that the equivalence provided by \eqref{eq:reduction-intro} is the same functor as in \eqref{eq:reduction-to-classical-intro}.

\subsection{The rigid refined local Langlands correspondence}

One immediate consequence of the extended spectral action constructed in this paper is a statement of the categorical local Langlands conjecture for arbitrary connected reductive groups. It is also likely that adapting the methods of \cite{GLCV} to the Fargues--Scholze setting, as explained in the previous subsection, gives another approach for producing this maximally general formulation of the CLLC. 

Returning temporarily to the non-categorical setting, we want to emphasize that the purpose of the rigid refined local Langlands conjecture is not just to formulate a parametrization of the $L$-packets for arbitrary connected reductive groups. Indeed, the first completely general version\footnote{By ``refined local Langlands conjectures'' in this context, we just mean a parametrization of each $L$-packet $\Pi_{\varphi}$ using some datum related to the centralizer of $\varphi$ in $\widehat{G}$.} of the refined local Langlands conjectures is already given in \cite{Arthur06} by taking the preimage of the $\widehat{G}_{\tn{ad}}$--image of $S_{\varphi} := Z_{\widehat{G}}(\varphi)$ in $\widehat{G}_{\tn{sc}}$ in order to parametrize the $L$-packets for all inner forms of a fixed quasi-split group. However, in the same paper, Arthur observes that this approach cannot yield normalized absolute transfer factors in general. 

By defining a version of the refined local Langlands correspondence parametrized by rigid inner twists, Kaletha proves:
\begin{thm}[\protect{\cite[Proposition 5.6]{Kal16}}]
    Given a tempered $L$-parameter $\varphi$ for a connected reductive group $G$ with endoscopic datum $\mathfrak{e}$, one can always construct a normalized absolute transfer factor $\Delta$ for $(G,\varphi,\mathfrak{e})$.
\end{thm}

Defining such transfer factors is the primary motivation for the aforementioned rigid formulation \cite{Kal16} of the refined local Langlands correspondence, which thus both encompasses all connected reductive groups and allows for a general statement of the endoscopic character identities.

One of the eventual aims of the categorical local Langlands project writ large is to bring the notion of endoscopy for $p$-adic groups under the umbrella of Fargues--Scholze. Since rigid inner twists are required to state the endoscopic character identities for groups which are not pure inner forms of a quasi-split group, it is reasonable to believe that one needs to use the extended version of the CLLC given in this paper in order to achieve this goal in full generality.

Furthermore, we hope that in the future one can use the extended geometrization to prove a version of the refined local Langlands correspondence (inspired by an earlier version of \cite{Fargues22}), which, among other desiderata, predicts:

\begin{conj}[\protect{\cite[Conjecture 5.1]{Dillery26}}] For a fixed $L$-parameter $\varphi$ for quasi-split $G$ with Whittaker datum $\mathfrak{w}$, there is a bijection
\begin{equation}\label{eq:rigLLC}
 \Pi_{\varphi}^{e}  \xrightarrow{\iota_{\mathfrak{w}}}  \tn{Irr}(S_{\varphi}^{+,\natural}),
\end{equation}
where $\Pi_{\varphi}^{e}$ is a $|\Bun_{G,\tn{basic}}^{e}|$-indexed union of the $L$-packets for $\varphi$ and the corresponding inner forms of $G$ and (denoting by $S_{\varphi}^{+}$ the preimage of $S_{\varphi}$ in $\widehat{G}^{e}$)
\begin{equation*}
S_{\varphi}^{+,\natural} := \frac{S_{\varphi}^{+}}{(S_{\varphi}^{+} \cap \widehat{G}_{\tn{sc}})^{\circ}}.
\end{equation*}
\end{conj}

We conclude by explaining one potential global application. As explained in the introduction to \cite{Kaletha18b}, rigid inner twists and (a variant of) the correspondence \eqref{eq:rigLLC} are required to state the conjectural multiplicity formula for discrete automorphic representations of a connected reductive group $G$, even when $Z_{G}$ is connected (more precisely, one needs $G$ to satisfy the Hasse principle and have connected center in order to use extended pure inner forms instead for this purpose). In the setting of a global function field one has the recent work \cite{DLH26} of Li-Huerta which reformulates the global Langlands conjecture in terms of the CLLC. We hope that synthesizing this paper and its techniques with \cite{DLH26} can make progress towards proving cases of the aforementioned multiplicity formula (cf. \cite[Conjecture 5.6]{Dillery26}) over global function fields.


\subsection{Notation and conventions}

\subsubsection{Fields}

We let $F$ be a non-archimedean local field with residue field $\Fq$ of characteristic $p$. We denote by $\breve{F}, \overline{F}$ the completion of the maximal unramified extension and an algebraic closure of $F$. We let $C$ be a fixed algebraically closed perfectoid field over $\Fqb$ together with an isomorphism $C^{\sharp} \xrightarrow{\sim} \hat{\ov{F}}$. 

We denote by $\Ocal_?$ the ring of integers of $? \in \{F, \breve{F}, \overline{F}, C\}$.

We shall denote by $\Gal_F, \Inert_F, P_F$ the absolute Galois group of $F$, its inertia subgroup and wild inertia respectively. We write $\Weil_F$ for the Weil group of $F$. 

\subsubsection{Geometry}

We denote by $\Perfd$ the category of perfectoid spaces over $\Fqb$, as in \cite[\S 7.1]{SW20} and by $\Perfd_C$ the category of perfectoid spaces over $\Spa(C)$. For $S \in \Perfd$ we denote by $X_{S}$ the associated relative Fargues--Fontaine curve, as in \cite[Definition II.1.15]{Geometrization}, and denote by $\Div^{1}$ the moduli stack over $\Perfd$ parametrizing degree $1$ relative Cartier divisors on $X_{S}$ (as in \cite[Definition II.1.19]{Geometrization}). For $C$ an algebraically closed perfectoid field, denote by $\Spd(C)$ the diamond associated to $\Spa(C,C^{+})$ (as in \cite[\S 8.3]{SW20}).

The datum of the isomorphism $C^{\sharp} \xrightarrow{\sim} \hat{\ov{F}}$ amounts to the choice of a point 
$$y : \Spa(C,\Ocal_C) \to \Div^1.$$
When the context is clear, we drop the $y$ and just write $\Spa(C) \to \Div^1$.


\subsubsection{Categories of sheaves}

We let $\ell \neq p$ be a prime and we let $\Lambda$ be a coefficient ring. When $\Lambda$ is torsion of residue characteristic $\ell$, for a diamond or $v$-stack $X$, we denote by 
$$\mc{D}(X, \Lambda)$$ 
the category of étale sheaves of $\Lambda$-modules on $X$. We shall write $\mc{D}(X)$ if $\Lambda$ is clear from the context. For $\ell$-adic coefficients (i.e. when $\Lambda = \Zlb, \Qlb$), we shall, as in \cite{Geometrization}, use the formalism of solid and lisse sheaves, see Section \ref{sec:Extended-Hecke-Action}. 

\subsubsection{Reductive groups}

We let $G$ be a quasi-split reductive group over $F$ together with a Borel pair $(B,T)$. We denote by $(X^*(T), X_*(T), \Phi, \Phi^{\vee}, \Delta, \Delta^{\vee})$, the based root datum of $(B,T)$ and by $W$ the Weyl group of $(B,T)$. We write $Z_G$ for the center of $G$.
We denote by $\widehat{G}$ the Langlands dual group, it is a canonically defined pinned reductive group over $\Lambda$ and by ${^L}G = \widehat{G} \rtimes \Gal_F$ with Galois action normalized so that for $\alpha \in \Delta^{\vee}$ (a simple root of $\widehat{G}$), we have a $\Gal_F$-equivariant isomorphism 
$$\widehat{G}_{\alpha} \cong \Ga(1),$$
see \cite[Section VI.11]{Geometrization}.

To work in the extended setting, we denote by $\widehat{G}^e$ the extended dual group as defined in \cite[\S 12.1]{Fargues22}. This is the (pro)-reductive group over $\Lambda$ given as
$$\widehat{G}^e = \varprojlim_Z \widehat{G/Z}$$
where the limit runs through the set of all finite central subgroups $Z \subset Z_{G}$.

\subsubsection{$\infty$-categories}

We shall write $\Pr$ for the category of presentable, stable, $\Lambda$-linear categories, equipped with its usual Lurie tensor product, we write $\Alg(\Pr)$ (resp. $\CAlg(\Pr)$) to denote the category of algebras in $\Pr$, this is the category of monoidal (resp. symmetric monoidal) categories in $\Pr$ whose tensor product commutes with colimits in both variables.

\subsubsection{Miscellaneous}
For $X$ a topological space and $A$ a topological abelian group, denote by $\mathcal{C}(X,A)$ the group of continuous functions from $X$ to $A$. When $A$ is an abstract abelian group and we do not specify the topology, then it is assumed that $A$ has the discrete topology.

\subsection{Acknowledgements}
The second author thanks the Max Planck Institute for Mathematics for its hospitality during the preparation of this paper. We thank Dori Bejleri, Patrick Brosnan, Tom Haines, David Hansen, Jessica Fintzen, Tasho Kaletha, Alex Youcis, and Konrad Zou for comments and discussions during the preparation of this paper. We thank Thibaud van den Hove for his very helpful clarifications about Hecke stacks. We also thank Yifei Zhao for explaining to us the content of \cite{GLCV} and discussions during the preparation of this paper. 


